\documentclass[10pt,a4paper]{amsart}
\usepackage{amsmath,amssymb,amsthm,graphicx}
\usepackage{verbatim}
\newenvironment{reference}{\comment}{\endcomment}
%\newenvironment{reference}{}{}
\newenvironment{slogan}{\comment}{\endcomment}
\newenvironment{history}{\comment}{\endcomment}

% For commutative diagrams we use Xy-pic
\usepackage[all]{xy}

% We use 2cell for 2-commutative diagrams.
\xyoption{2cell}
% Theorem environments.
%
\theoremstyle{plain}
\newtheorem{theorem}[subsection]{Theorem}
\newtheorem{proposition}[subsection]{Proposition}
\newtheorem{lemma}[subsection]{Lemma}

\theoremstyle{definition}
\newtheorem{definition}[subsection]{Definition}
\newtheorem{example}[subsection]{Example}
\newtheorem{exercise}[subsection]{Exercise}
\newtheorem{situation}[subsection]{Situation}

\theoremstyle{remark}
\newtheorem{remark}[subsection]{Remark}
\newtheorem{remarks}[subsection]{Remarks}

\numberwithin{equation}{subsection}

% Macros
%
\def\lim{\mathop{\mathrm{lim}}\nolimits}
\def\colim{\mathop{\mathrm{colim}}\nolimits}
\def\Spec{\mathop{\mathrm{Spec}}}
\def\Hom{\mathop{\mathrm{Hom}}\nolimits}
\def\Ext{\mathop{\mathrm{Ext}}\nolimits}
\def\SheafHom{\mathop{\mathcal{H}\!\mathit{om}}\nolimits}
\def\SheafExt{\mathop{\mathcal{E}\!\mathit{xt}}\nolimits}
\def\Sch{\mathit{Sch}}
\def\Mor{\mathop{\mathrm{Mor}}\nolimits}
\def\Ob{\mathop{\mathrm{Ob}}\nolimits}
\def\Sh{\mathop{\mathit{Sh}}\nolimits}
\def\NL{\mathop{N\!L}\nolimits}
\def\CH{\mathop{\mathrm{CH}}\nolimits}
\def\proetale{{pro\text{-}\acute{e}tale}}
\def\etale{{\acute{e}tale}}
\def\QCoh{\mathit{QCoh}}
\def\Ker{\mathop{\mathrm{Ker}}}
\def\Im{\mathop{\mathrm{Im}}}
\def\Coker{\mathop{\mathrm{Coker}}}
\def\Coim{\mathop{\mathrm{Coim}}}

% Boxtimes
%
\DeclareMathSymbol{\boxtimes}{\mathbin}{AMSa}{"02}

%
% Macros for moduli stacks/spaces
%
\def\QCohstack{\mathcal{QC}\!\mathit{oh}}
\def\Cohstack{\mathcal{C}\!\mathit{oh}}
\def\Spacesstack{\mathcal{S}\!\mathit{paces}}
\def\Quotfunctor{\mathrm{Quot}}
\def\Hilbfunctor{\mathrm{Hilb}}
\def\Curvesstack{\mathcal{C}\!\mathit{urves}}
\def\Polarizedstack{\mathcal{P}\!\mathit{olarized}}
\def\Complexesstack{\mathcal{C}\!\mathit{omplexes}}
% \Pic is the operator that assigns to X its picard group, usage \Pic(X)
% \Picardstack_{X/B} denotes the Picard stack of X over B
% \Picardfunctor_{X/B} denotes the Picard functor of X over B
\def\Pic{\mathop{\mathrm{Pic}}\nolimits}
\def\Picardstack{\mathcal{P}\!\mathit{ic}}
\def\Picardfunctor{\mathrm{Pic}}
\def\Deformationcategory{\mathcal{D}\!\mathit{ef}}
\usepackage{cleveref}
\crefname{originalreference}{original reference}{original references}
\Crefname{originalreference}{Original reference}{Original references}
\makeatletter
\newlabel{lemma-bound-dualizing-implies-bound}{{0B4S}{source}{}{sourcepointer-1}{}}
\newlabel{lemma-bound-dualizing-implies-bound@cref}{{[originalreference][0][]0B4S}{[1][0][]source}{}{}{}}
\makeatother
\crefname{originalreference}{original reference}{original references}
\Crefname{originalreference}{Original reference}{Original references}
\makeatletter
\newlabel{proposition-Grauert-Riemenschneider}{{0AXD}{source}{}{sourcepointer-2}{}}
\newlabel{proposition-Grauert-Riemenschneider@cref}{{[originalreference][0][]0AXD}{[1][0][]source}{}{}{}}
\makeatother
\crefname{originalreference}{original reference}{original references}
\Crefname{originalreference}{Original reference}{Original references}
\makeatletter
\newlabel{situation-rational}{{0B4W}{source}{}{sourcepointer-3}{}}
\newlabel{situation-rational@cref}{{[originalreference][0][]0B4W}{[1][0][]source}{}{}{}}
\makeatother
\title{The trace map from the dualizing module of the resolution to the dualizing module of the base}
\author{The Stacks Project Authors}
\date{stacks-project:resolve.tex:local-snapshot-20261012; GNU FDL 1.2 or later}
\begin{document}
\maketitle
\noindent\textit{Editorial note.} All delivered mathematical text is the author's own, copied from the pinned Stacks Project chapter file resolve.tex; no statement or proof step was rewritten, reordered or removed.\par\medskip
\noindent\textit{Editorial note.} The counted proof is the author's complete argument: the shortening of the duality sequence through the Grauert-Riemenschneider vanishing, and the vanishing of the last term through the definition of a rational singularity.\par\medskip
\noindent\textit{Editorial note.} The label lemma-dualizing-rational names the counted statement, and the proof environment immediately following it in the source is delivered verbatim as its proof.\par\medskip
\noindent\textit{Editorial note.} The counted statement is carried by the author's own lemma environment in the chapter source.\par\medskip
\noindent\textit{Editorial note.} Pre-open gate: python operations/agent-authoring/tools/precard\_gate.py --source resolve.tex --label lemma-dualizing-rational --start 2617 --end 2638 --local-id 28121 returned RESULT=OK (L28121 free, occupancy FREE, public interval clear). Pre-assembly tag self-check: canonical tag resolved and not present in reference/public-items.jsonl.\par\medskip
\noindent\textit{Editorial note.} The statement and proof are copied from the pinned chapter file; the cited inputs are kept as external pointers with their canonical Stacks tags.\par\medskip
\noindent\textit{Allowed standard prerequisites (editorial).}\begin{itemize}
\item Vanishing of the first direct image of the dualizing module in the rational situation. Same chapter, the Grauert-Riemenschneider vanishing statement for the first direct image of the dualizing module. Identifies the derived pushforward of the dualizing module with its zeroth pushforward, which is what makes the duality sequence short.
\item The duality sequence relating the pushforward of the dualizing module to the dualizing complex of the base. Same chapter, the argument accompanying the lemma that a bound on the dualizing sheaf implies a bound, containing the displayed short exact sequence. Provides the short exact sequence whose last term is an Ext group of the first direct image of the structure sheaf.
\item Vanishing of the first cohomology of the structure sheaf for a rational singularity. The definition of a rational singularity, giving the vanishing of the first cohomology of the structure sheaf for every modification. Kills the last term of the short exact sequence.
\end{itemize}\medskip
\ifdefined\hypertarget\hypertarget{sourcepointer-1}{}\fi
\noindent\textit{Source reference.} Original reference 0B4S: sources/stacks-resolve/pinned-original-source.tex, source lines 2233--2233. Prerequisite pointer invoked at proof line(s) 2635; it names a standard fact of the same chapter that the counted argument uses, and every step of the counted proof built on it is retained verbatim in the delivered body.\par\medskip
\ifdefined\hypertarget\hypertarget{sourcepointer-2}{}\fi
\noindent\textit{Source reference.} Original reference 0AXD: sources/stacks-resolve/pinned-original-source.tex, source lines 1838--1838. Prerequisite pointer invoked at proof line(s) 2628; it names a standard fact of the same chapter that the counted argument uses, and every step of the counted proof built on it is retained verbatim in the delivered body.\par\medskip
\ifdefined\hypertarget\hypertarget{sourcepointer-3}{}\fi
\noindent\textit{Source reference.} Original reference 0B4W: sources/stacks-resolve/pinned-original-source.tex, source lines 2455--2455. Prerequisite pointer invoked at proof line(s) 2619; it names a standard fact of the same chapter that the counted argument uses, and every step of the counted proof built on it is retained verbatim in the delivered body.\par\medskip
\begin{lemma}
\label{lemma-dualizing-rational}
In Situation \ref{situation-rational} assume $A$ has a
dualizing complex $\omega_A^\bullet$. With $\omega_X$ the dualizing
module of $X$, the trace map $H^0(X, \omega_X) \to \omega_A$ is an
isomorphism and consequently there is a canonical map
$f^*\omega_A \to \omega_X$.
\end{lemma}

\begin{proof}
By Grauert-Riemenschneider
(Proposition \ref{proposition-Grauert-Riemenschneider}) we see that
$Rf_*\omega_X = f_*\omega_X$. By duality we have a short exact
sequence
$$
0 \to f_*\omega_X \to \omega_A \to
\Ext^2_A(R^1f_*\mathcal{O}_X, \omega_A) \to 0
$$
(for example see proof of Lemma \ref{lemma-bound-dualizing-implies-bound})
and since $A$ defines a rational singularity we obtain
$f_*\omega_X = \omega_A$.
\end{proof}

\end{document}
